% ============================================================
% CHAPTER 2: SET COVER (matches Vazirani Ch. 2)
% ============================================================
\setcounter{chapter}{1}
\chapter{Set Cover}
\label{ch:setcover}

\section{Intuition: The Greedy Instinct}

\begin{intuitionbox}
\textbf{The Greedy Idea:} When you need to cover all elements, pick the set that gives you the most ``bang for your buck'' — the most new elements per unit cost. Repeat until everything is covered.

\textbf{Why it works:} The price paid for each element when it's covered serves as a \emph{dual variable}. Summing all prices gives a lower bound on OPT (via LP duality, Chapter 13). The harmonic number $H_n$ emerges naturally because early elements get covered cheaply, later elements pay more.
\end{intuitionbox}

\section{Problem Definition}

\begin{definition}[Set Cover]
Given a universe $U = \{e_1, \dots, e_n\}$ and a collection of subsets $\mathcal{S} = \{S_1, \dots, S_m\}$ with costs $c(S_i) \geq 0$, find a minimum-cost subcollection $\mathcal{C} \subseteq \mathcal{S}$ such that $\bigcup_{S \in \mathcal{C}} S = U$.
\end{definition}

\section{The Greedy Algorithm}

\begin{algorithmdef}[Greedy Set Cover — $H_n$ Approximation]
\begin{enumerate}
    \item $C \gets \emptyset$ (covered elements)
    \item While $C \neq U$:
    \begin{enumerate}
        \item Pick $S \in \mathcal{S}$ minimizing $\frac{c(S)}{|S \setminus C|}$ (cost-effectiveness)
        \item For each $e \in S \setminus C$, set $\text{price}(e) = \frac{c(S)}{|S \setminus C|}$
        \item $C \gets C \cup S$
    \end{enumerate}
    \item Output the picked sets
\end{enumerate}
\end{algorithmdef}

\begin{theorem}
The Greedy Set Cover algorithm achieves approximation factor $H_n = 1 + \frac{1}{2} + \dots + \frac{1}{n} \approx \ln n + \gamma$.
\end{theorem}
\begin{proof}[Proof Sketch]
Each element's price is a lower bound on its contribution to OPT. The $k$-th element covered pays at most $\frac{\OPT}{k}$. Summing gives $H_n \cdot \OPT$.
\end{proof}

\subsection{Layering: Another 2-Approximation for Weighted Vertex Cover}

Layering decomposes vertex weights into degree-weighted functions. Each layer gives a factor-2 vertex cover, and summing gives factor-2 for the weighted case. This is a preview of the primal-dual schema (Chapter 22).

\subsection{Application to Shortest Superstring}

The Shortest Superstring problem (find shortest string containing all given strings as substrings) reduces to Set Cover via string overlaps. Greedy gives $2H_n$-approximation; Li~\cite{li92} improved to factor 3.

\section{Other Algorithms for Set Cover}

\begin{table}[H]
\centering
\caption{Approaches to Set Cover}
\begin{tabular}{lccc}
\toprule
Algorithm & Approximation & Technique \\
\midrule
Greedy (this chapter) & $H_n$ & Combinatorial \\
LP Rounding (Ch. 14) & $f$ (max frequency) & LP + Rounding \\
Primal-Dual (Ch. 13, 22) & $f$ & Primal-Dual Schema \\
Randomized Rounding & $O(\log n)$ & LP + Randomization \\
\bottomrule
\end{tabular}
\end{table}

\textbf{Why Greedy?} It's simple, fast, and achieves the optimal $\ln n$ bound (Feige~\cite{feige98} proved no $(1-\epsilon)\ln n$ unless P = NP).

\textbf{Why not better?} Feige~\cite{feige98} showed no polynomial-time algorithm can achieve $(1-\epsilon)\ln n$ approximation unless NP $\subseteq$ DTIME$(n^{\log\log n})$.

\section{Python Implementation}

\begin{lstlisting}[style=python, caption=Greedy Set Cover]
def greedy_set_cover(universe, sets, costs):
    """H_n-approximation for Set Cover."""
    covered = set()
    picked = []
    total_cost = 0.0
    prices = {}
    
    while covered != universe:
        best_set = None
        best_ratio = float('inf')
        for s_id, s in sets.items():
            new = s - covered
            if not new: continue
            ratio = costs[s_id] / len(new)
            if ratio < best_ratio:
                best_ratio = ratio
                best_set = s_id
        
        if best_set is None: break
        new_elements = sets[best_set] - covered
        for e in new_elements:
            prices[e] = best_ratio
        covered |= sets[best_set]
        picked.append(best_set)
        total_cost += costs[best_set]
    
    return picked, total_cost
\end{lstlisting}

\section{Applications}
\begin{itemize}
    \item \textbf{Feature selection:} Minimum features covering all classes
    \item \textbf{Antenna placement:} Cover all locations with minimum transmitters
    \item \textbf{Test case minimization:} Cover all code paths with fewest tests
    \item \textbf{Emergency facilities:} Fire stations covering all districts
    \item \textbf{Data compression:} Shortest superstring via set cover
\end{itemize}

\subsection*{Real Example: A/B Test Minimization at a Tech Company}
\textbf{Scenario:} A/B testing platform runs 10,000 experiments. Each experiment tests a variant against control on a subset of 100 metrics (click-through, revenue, retention, etc.). Running an experiment on all metrics is expensive.
\textbf{Model:} Universe = 100 metrics. Sets = experiments (each covers $\sim$15 metrics). Cost = compute cost per experiment.
\textbf{Why greedy:} The greedy algorithm picks the experiment covering the most \emph{new} metrics per unit cost. $H_{100} \approx 5.18$ means greedy uses $\leq 5\times$ the minimum experiments.
\textbf{Real result:} In practice, greedy reduces experiments by $70\%$ vs. running all, with $H_n$ bound giving stakeholders confidence.
\textbf{Set cover variant:} \emph{Budgeted maximum coverage} — given $k$ experiments, maximize metrics covered. Greedy gives $(1-1/e)$-approx (Chapter 13).

\subsection*{Real Example: Feature Selection for Medical Diagnosis}
\textbf{Scenario:} A hospital has $n=10,000$ patient features (lab results, vitals, history). They want to select the minimum subset that still predicts $m=50$ diseases with $95\%$ accuracy.
\textbf{Model:} Universe $U$ = diseases. Each feature $f$ covers the subset $S_f \subseteq U$ it helps predict above threshold. Cost $c(f) = 1$ (each feature has equal acquisition cost).
\textbf{Why greedy works here:} At each step, greedy picks the feature that adds the most \emph{new} diseases predicted. The $H_{50} \approx 4.5$ guarantee means if optimal uses 20 features, greedy uses $\leq 90$. In practice, greedy often finds near-optimal because features have diminishing returns (submodularity).
\textbf{Real constraint:} Some features are expensive (MRI) vs. cheap (blood pressure). Weighted set cover handles this — greedy picks by cost-effectiveness.
\textbf{Alternative:} LASSO (convex relaxation) gives exact recovery under restricted isometry but no worst-case guarantee; ILP is exact but slow for $n=10,000$.

\section{Tight Example}

\begin{example}
Universe $\{e_1, \dots, e_n\}$. Sets: $S_0 = U$ (cost $1+\epsilon$), $S_i = \{e_i\}$ (cost $1/i$). Greedy picks all $S_i$ (total $H_n$). Optimal picks $S_0$ (cost $1+\epsilon$). Ratio $\to H_n$.
\end{example}

\section{Exercises}
\begin{exercise}
Show greedy achieves $H_k$ where $k = \max_i |S_i|$.
\end{exercise}
\begin{exercise}
Give a factor-$f$ algorithm where $f$ is maximum element frequency.
\end{exercise}

